% !TEX root = ../my-thesis.tex

\chapter{Diskussion}
\label{sec:diskussion}

Dieses Kapitel ordnet die Ergebnisse beider Untersuchungsebenen
ein. Zunächst werden die vier Forschungsfragen beantwortet,
anschließend werden die Grenzen der Arbeit dargelegt.

\section{Beantwortung der Forschungsfragen}
\label{sec:diskussion:fragen}

\paragraph{F1: Modellierung der beiden Workflow-Muster.}
Sowohl das Pipeline- als auch das Scatter-Gather-Muster ließen sich
in allen fünf Systemen mit denselben Benchmark-Skripten umsetzen.
Für die beiden untersuchten Muster stellt die grundsätzliche
Abbildbarkeit damit kein wesentliches Unterscheidungsmerkmal der
Systeme dar. Unterschiede zeigen sich vielmehr in der Art der
Workflow-Beschreibung und darin, wie die Beziehungen zwischen den
einzelnen Schritten ausgedrückt werden. Die
beobachteten Beschreibungsformen reichen von systemeigenen Formaten
über den Standard CWL bis zu Programmierschnittstellen und lassen
sich damit in die von Suter et al. unterschiedenen Kategorien
einordnen~\cite[S.~8]{suter2026terminology}.

Beim Scatter-Gather-Muster treten diese Unterschiede deutlicher
hervor. Neben den Abhängigkeiten müssen dort die Aufteilung in
parallele Verarbeitungsinstanzen und deren anschließende
Zusammenführung ausgedrückt werden. Auch eine Variation der
Teilstückzahl erfordert je nach System unterschiedliche Änderungen
an der Workflow-Beschreibung. Die Umsetzung beider Muster zeigt
damit, dass sich die betrachteten Systeme bei diesen beiden Mustern
weniger durch ihre Ausdrucksfähigkeit als durch die Art der
Workflow-Beschreibung und den Aufwand bei Änderungen unterscheiden.

\paragraph{F2: Laufzeitverhalten der Workflow-Management-Systeme.}
Auf beiden Untersuchungsebenen zeigt sich ein ähnliches Bild beim
Overhead. StreamFlow, Nextflow und Merlin liegen in einem
gemeinsamen Bereich, AiiDA darüber, und Pegasus in der lokalen
Untersuchung mit deutlichem Abstand an der Spitze. Dass
diese Anordnung auf einem einzelnen Rechner und auf einem
dedizierten Rechenknoten gleichermaßen auftritt, spricht dagegen,
sie allein der jeweiligen Ausführungsumgebung zuzuschreiben. Ein
proportionaler Zusammenhang zwischen Rechenlast und Overhead ist in
keinem der Systeme erkennbar, weshalb der Overhead mit wachsender
Rechenlast gegenüber der Rechenzeit an Gewicht verliert. Dies
gehört zu den Gründen, aus denen die Aufteilung auf mehrere
Teilstücke beim Scatter-Gather-Muster erst ab der mittleren
Rechenlast spürbare Gewinne bringt, während bei der kurzen
Rechenlast der Overhead einen größeren Teil der Gesamtlaufzeit ausmacht.

Die beiden Koordinationsgrößen machen sichtbar, an welchen Stellen zusätzliche Zeit anfällt. Besonders deutlich zeigt sich dies bei Pegasus und AiiDA. Bei Pegasus fällt die mediane Übergangslatenz von rund 20 Sekunden mit dem höchsten gemessenen Overhead zusammen. Bei AiiDA fallen die Übergangslatenzen mit ein bis zwei Sekunden geringer aus, liegen aber auf beiden Ebenen deutlich über denen der drei übrigen Systeme und erklären bereits im Pipeline-Muster einen wesentlichen Teil des Overheads. Die Ergebnisse sprechen damit dafür, dass die höheren Overhead-Werte von Pegasus und AiiDA zu einem erheblichen Teil mit der Koordination und Übergabe
zwischen aufeinanderfolgenden Workflow-Schritten zusammenhängen.

Die Startspreizung ist davon getrennt zu betrachten. Sie beschreibt, wie schnell ein System die parallelen Instanzen startet. Merlin zeigt das deutlich: Lokal liegt seine
Startspreizung bei vier Teilstücken mit rund 0,73 Sekunden in
derselben Größenordnung wie bei AiiDA, während sein Overhead
deutlich niedriger bleibt. Eine größere Startspreizung führt damit
nicht zwangsläufig zu einem höheren Overhead.

\paragraph{F3: Betrieb auf dem HPC-System.}
Vier der fünf untersuchten Systeme ließen sich auf MOGON unter den
Rechten eines gewöhnlichen Nutzerkontos bereitstellen und
betreiben, Pegasus konnte nicht einbezogen werden. Dabei
unterscheiden sich nicht nur die erforderlichen
Einrichtungsschritte, sondern auch die Betriebsmodelle unter Slurm.
Nextflow, StreamFlow und AiiDA reichen die Workflow-Schritte
einzeln als Slurm-Jobs ein. Die Steuerung erfolgt bei Nextflow und
StreamFlow durch einen Vordergrundprozess und bei AiiDA durch einen
dauerhaft laufenden Daemon. Merlin verwendet dagegen eine
gemeinsame Worker-Allokation, innerhalb derer die einzelnen
Workflow-Schritte abgearbeitet werden. Damit fallen Warteschlange
und Startphase eines Slurm-Jobs bei den ersten drei Systemen für
jeden eingereichten Schritt erneut an, bei Merlin dagegen nur
einmal beim Start der Allokation.

Für die Laufzeitzusammensetzung zeigt der reguläre Slurm-Betrieb,
dass die Ausführungsumgebung einen wesentlich größeren Einfluss
haben kann als das Workflow-Management-System selbst. In den
Beobachtungsläufen schwankte die Wartezeit in der
Slurm-Warteschlange zwischen 6 und 17 Sekunden je Job. Hinzu kamen
vor Beginn der eigentlichen Nutzlast durchgehend mehr als
40 Sekunden. Die anschließende systemspezifische Untersuchung
zeigte dagegen, dass der innerhalb dieses Zeitraums
identifizierbare Anteil der Workflow-Management-Systeme bei den
Systemen mit einzelnen Slurm-Jobs lediglich zwischen 55 und
104 Millisekunden lag. Eine Messung im regulären Clusterbetrieb würde damit überwiegend Wartezeit und die Zeit vor Beginn der Nutzlast erfassen und die Unterschiede zwischen den Systemen überlagern. Für den eigentlichen Systemvergleich wurde deshalb eine bereits bestehende Slurm-Allokation verwendet, innerhalb derer
Warteschlange und Jobstart nicht erneut für jeden Workflow-Schritt
anfielen.

Für den HPC-Betrieb ist damit auch die Einbindung in die
Ausführungsumgebung von Bedeutung. Die Art der Slurm-Einbindung beeinflusst, wie häufig die vom Cluster verursachten Zeitanteile innerhalb eines
Workflows anfallen. Bei AiiDA kam mit dem Ablageort von Datenbank
und Repository ein weiterer Einflussfaktor hinzu, dessen
Veränderung den gemessenen Overhead deutlich beeinflusste. Die gemessenen Laufzeiten hängen auf dem HPC-System daher nicht nur
vom Workflow-Management-System selbst, sondern auch von den
Bedingungen der Ausführungsumgebung ab.

\paragraph{F4: Eignung für Workflow-Arten und Anwendungsfälle.}
Aus den Ergebnissen ergibt sich keine allgemeine Rangfolge der
Systeme. Welches System geeignet erscheint, hängt vielmehr davon
ab, welche Eigenschaften für den jeweiligen Anwendungsfall im
Vordergrund stehen.

Nextflow beschreibt Workflows datenflussorientiert und stellt die
schnelle Entwicklung, parallele Ausführung und Reproduzierbarkeit
in den Mittelpunkt~\cite{ditommaso2017nextflow}. Beide untersuchten
Muster ließen sich kompakt über Prozesse und Channels abbilden, wobei
auch Änderungen der Teilstückzahl keine Anpassung der
Workflow-Struktur erforderten. Zusammen mit dem geringen Overhead,
den niedrigen Koordinationszeiten und der direkten Slurm-Anbindung
ergibt sich damit ein günstiges Profil für datenflussorientierte
Workflows mit sequenziellen und parallelen Verarbeitungsschritten.
Im regulären Clusterbetrieb ist allerdings zu berücksichtigen, dass
jeder Prozess als eigener Slurm-Job eingereicht wird und damit die
Warte- und Startzeiten des Ressourcenverwalters je Schritt erneut
entstehen können.

StreamFlow ist auf Workflows in hybriden Ausführungsumgebungen
ausgerichtet und trennt die Workflow-Beschreibung von der
Konfiguration der
Ausführungsumgebung~\cite{colonnelli2021streamflow}. Diese Trennung
zeigte sich beim Wechsel auf MOGON: Die in CWL formulierte
Workflow-Beschreibung blieb unverändert, während die
Ausführungskonfiguration angepasst wurde. Zusammen mit dem geringen Overhead und den niedrigen gemessenen Koordinationszeiten bietet sich StreamFlow
damit insbesondere an, wenn eine standardisierte
Workflow-Beschreibung und die Trennung von Workflow und
Ausführungsumgebung im Vordergrund stehen. Einschränkend war die
Bereitstellung der untersuchten Vorabversion aufwendiger und
erforderte Anpassungen für die Slurm-Ausführung.

\begin{sloppypar}
Merlin ist für große Ensembles von Rechenaufgaben auf HPC-Systemen
ausgelegt und verteilt Aufgaben über ein
Erzeuger-Verbraucher-Modell an
Worker~\cite{peterson2022merlin}. Dieses Modell ließ sich beim
Scatter-Gather-Muster für die parallelen Verarbeitungsinstanzen
nutzen. Auf MOGON wurden die Workflow-Schritte innerhalb einer
gemeinsamen Worker-Allokation ausgeführt. Dadurch fielen die
Warte- und Startzeiten des Ressourcenverwalters nicht für jeden
Schritt erneut an. Merlin bietet sich damit insbesondere an, wenn viele Rechenaufgaben innerhalb einer einmal angeforderten Allokation abgearbeitet werden
sollen, ohne dass jede Aufgabe erneut die Warteschlange durchläuft. Die von Merlin
angestrebte Skalierung auf sehr große Aufgabenzahlen wurde in
dieser Arbeit jedoch nicht untersucht, da das
Scatter-Gather-Muster höchstens vier parallele
Verarbeitungsinstanzen umfasste.
\end{sloppypar}

AiiDA stellt die dauerhafte Nachvollziehbarkeit wissenschaftlicher
Abläufe in den Mittelpunkt und hält Prozesse und Daten während der
Ausführung in einem Provenance-Graphen
fest~\cite{huber2020aiida}. Dieser Schwerpunkt geht mit einem
umfangreicheren Betriebsmodell sowie in der Untersuchung mit
höheren Koordinationszeiten und höherem Overhead als bei Nextflow,
StreamFlow und Merlin einher. Zudem zeigte die Konfigurationsstudie
einen deutlichen Einfluss des Ablageorts der Datenhaltung auf die
gemessenen Zeiten. Damit stellt AiiDA insbesondere dann eine relevante Option dar, wenn die dauerhafte Nachvollziehbarkeit von Prozessen und Daten
für den Anwendungsfall stärker gewichtet wird als der in dieser
Untersuchung beobachtete höhere Overhead und der
zusätzliche Koordinations- und Betriebsaufwand, wobei die absolute Höhe der
gemessenen Zeiten von der Konfiguration der Datenhaltung abhängt.

Pegasus richtet sich auf großskalige, mehrstufige Simulations- und
Analyse-Abläufe auf verteilten
Rechenressourcen~\cite[S.~17]{deelman2015pegasus} und übersetzt eine
abstrakte Workflow-Beschreibung in einen ausführbaren
Workflow~\cite[S.~18]{deelman2015pegasus}. In der lokalen
Untersuchung wies Pegasus allerdings den mit Abstand höchsten
Overhead auf, der bei kurzen Rechenschritten besonders stark ins
Gewicht fiel. Die Dokumentation beschreibt den Aufwand beim
Einplanen kurzer Jobs selbst als Problemfall und sieht dafür das
Bündeln mehrerer Aufgaben vor~\cite{pegasusdocs-optimizing}.
Pegasus erscheint damit für die hier untersuchten feingranularen
Workflows weniger günstig, wenn ein geringer Laufzeit-Overhead im
Vordergrund steht. Die Eignung für das untersuchte HPC-System
konnte nicht bewertet werden, da sich Pegasus dort nicht
bereitstellen ließ.

Systemübergreifend zeigt sich, dass die Bedeutung des
Laufzeit-Overheads von der Größe der einzelnen Rechenaufgaben
abhängt. Bei kurzen Rechenschritten kann die zusätzlich durch das
Workflow-Management-System beanspruchte Zeit einen erheblichen
Anteil der Gesamtlaufzeit ausmachen. Mit zunehmender Rechenlast
verliert dieser Anteil dagegen an Gewicht. Für die Systemwahl können dann andere Eigenschaften stärker in den Vordergrund rücken, etwa die Art der Workflow-Beschreibung oder die
Nachvollziehbarkeit von Prozessen und Daten. Auf einem HPC-System
kommt hinzu, wie das jeweilige System den Ressourcenverwalter
einbindet. Im regulären Slurm-Betrieb fielen bei einzeln
eingereichten Jobs Warte- und Startzeiten für jeden
Workflow-Schritt erneut an, während Merlin die Schritte innerhalb
einer gemeinsamen Worker-Allokation ausführte. Welche dieser
Eigenschaften für die Eignung eines Systems entscheidend sind,
hängt damit von den Anforderungen des jeweiligen Anwendungsfalls
ab. Die Grenzen dieser Aussagen betrachtet
Abschnitt~\ref{sec:diskussion:grenzen}.

\section{Grenzen der Arbeit}
\label{sec:diskussion:grenzen}

Die Untersuchung betrachtet mit dem Pipeline- und dem
Scatter-Gather-Muster zwei grundlegende Workflow-Strukturen. Die
Ergebnisse zeigen, wie die fünf Systeme diese Muster modellieren
und ausführen. Komplexere Strukturen wie größere
Abhängigkeitsgraphen oder verschachtelte Workflows wurden nicht untersucht. Aussagen zur Modellierung und Eignung gelten daher für die beiden betrachteten
Muster.

Die Rechenschritte selbst bestehen aus bewusst einfach gehaltenen
Python-Skripten, damit alle Systeme dieselbe kontrollierte
Rechenlogik ausführen. Die verwendeten Rechenaufgaben bilden jedoch
keine reale wissenschaftliche Anwendung mit großen Datenmengen oder
intensiven Dateizugriffen ab. Inwieweit sich die beobachteten
Unterschiede unter solchen Anwendungslasten verändern, wurde daher
nicht untersucht.

Das Scatter-Gather-Muster wurde mit einem, zwei und vier
Teilstücken ausgeführt. Die Ergebnisse beschreiben damit das
Verhalten bei wenigen parallelen Verarbeitungsinstanzen. Ob sich
die beobachteten Unterschiede bei deutlich größeren Aufgabenzahlen
in gleicher Weise fortsetzen, lässt sich daraus nicht ableiten.
Dies betrifft insbesondere Merlin, dessen Entwurf ausdrücklich auf
große Ensembles zielt.

Die Messungen fanden auf einem einzelnen Rechner und auf einem
einzelnen HPC-System statt. Die absoluten Laufzeiten hängen damit
von diesen beiden Umgebungen ab und lassen sich nicht unmittelbar auf andere
Rechner oder HPC-Umgebungen übertragen.
Die innerhalb einer Untersuchungsebene beobachteten Unterschiede 
beziehen sich somit auf die jeweils verwendete Ausführungsumgebung.
Dies gilt auch für die systemspezifischen Befunde auf MOGON, 
beispielsweise die bei StreamFlow erforderliche Anpassung oder den bei AiiDA beobachteten Einfluss des Ablageorts der Datenhaltung. Ob diese Befunde unter
anderen HPC-Umgebungen in gleicher Weise auftreten, wurde nicht
untersucht.

Eine weitere Eingrenzung ergibt sich aus der Definition des
Overheads. Er wurde als zusätzlicher Zeitbedarf gegenüber einer
Referenzausführung derselben Rechenlogik ohne
Workflow-Management-System bestimmt. Die Messgröße gibt somit an,
wie viel Zeit zusätzlich anfällt, nicht, an welcher Stelle im
System sie entsteht. Die Übergangslatenzen und die Startspreizung
grenzen einzelne Anteile davon näher ein, erfassen jedoch nicht sämtliche internen Vorgänge eines Systems.

Jede Konfiguration wurde bei kurzer und mittlerer Rechenlast
fünfmal, bei langer Rechenlast dreimal ausgeführt und über den
Median ausgewertet. Die Verwendung des Medians reduziert den
Einfluss einzelner abweichender Läufe. Mit mehr Wiederholungen
ließe sich zusätzlich angeben, wie stark die Werte zwischen den
Läufen einer Konfiguration streuen.

Pegasus konnte auf MOGON nicht ausgeführt werden, da es sich
nicht bereitstellen ließ. Für dieses System liegen daher keine
Laufzeitdaten auf der HPC-Ebene vor, die quantitativen Aussagen zu
Pegasus stützen sich allein auf die lokale Untersuchung. Die
fehlende Bereitstellbarkeit von Pegasus unter den gegebenen
Bedingungen stellt dabei selbst ein Ergebnis der Untersuchung dar.